\documentclass[11pt,a4paper]{article}
\usepackage[utf8]{inputenc}
\usepackage{amsmath}
\usepackage{amssymb}
\usepackage{graphicx}
\usepackage{geometry}
\usepackage{float}
\usepackage{hyperref}
\usepackage{listings}
\usepackage{xcolor}

\geometry{margin=1in}

\title{Lane-Tracking Assist (LTA) Control System\\
\large Linear State-Feedback Control with Speed Management}
\author{Robotics Course Project}
\date{\today}

\begin{document}

\maketitle

\begin{abstract}
This document presents a comprehensive Lane-Tracking Assist (LTA) control system implementing linear state-feedback control from robotics fundamentals. The system uses sensor-based state estimation (GNSS, IMU, wheel encoders) and camera-based lane detection to provide autonomous lane-keeping and speed control capabilities. We demonstrate the control architecture, vehicle dynamics model, and experimental effectiveness of the approach.
\end{abstract}

\section{Introduction}

The Lane-Tracking Assist system implements a hierarchical control architecture with three main layers:
\begin{enumerate}
    \item \textbf{Perception Layer}: Sensor fusion and lane detection
    \item \textbf{Control Layer}: Linear feedback steering control and speed management
    \item \textbf{Actuation Layer}: Vehicle dynamics and physics simulation
\end{enumerate}

The system operates in three modes: \texttt{MANUAL} (no assistance), \texttt{WARNING} (visual warnings only), and \texttt{ASSIST} (active control intervention).

\section{Vehicle Dynamics Model}

\subsection{Kinematic Model}

The vehicle follows a bicycle model with Ackermann steering geometry. The state equations are:

\begin{align}
\dot{x} &= v \cos(\theta) \\
\dot{y} &= v \sin(\theta) \\
\dot{\theta} &= \frac{v \tan(\delta)}{L}
\end{align}

where:
\begin{itemize}
    \item $(x, y)$ is the vehicle position in world coordinates
    \item $\theta$ is the vehicle heading angle
    \item $v$ is the forward velocity
    \item $\delta$ is the steering angle
    \item $L$ is the wheelbase (distance between front and rear axles)
\end{itemize}

\subsection{Dynamic Forces}

The vehicle experiences several forces affecting its motion:

\paragraph{Longitudinal Dynamics:}
\begin{equation}
m \dot{v} = F_{\text{drive}} - F_{\text{drag}} - F_{\text{rolling}} - F_{\text{brake}}
\end{equation}

where:
\begin{align}
F_{\text{drive}} &= \tau \cdot \text{throttle} \quad \text{(engine force)} \\
F_{\text{drag}} &= \frac{1}{2} \rho C_d A v^2 \quad \text{(air resistance)} \\
F_{\text{rolling}} &= C_r m g \quad \text{(rolling resistance)} \\
F_{\text{brake}} &= \beta \cdot \text{brake} \quad \text{(braking force)}
\end{align}

\paragraph{Lateral Dynamics:}
Lateral acceleration during cornering is constrained by tire friction:
\begin{equation}
a_{\text{lat}} = \frac{v^2}{R} \leq \mu g
\end{equation}

where $R$ is the curve radius and $\mu$ is the coefficient of friction.

\subsection{Implementation Parameters}

The simulator uses realistic parameters:
\begin{itemize}
    \item Wheelbase: $L = 2.7$ m (typical sedan)
    \item Mass: $m = 1500$ kg
    \item Drag coefficient: $C_d = 0.3$
    \item Frontal area: $A = 2.2$ m$^2$
    \item Maximum steering angle: $\delta_{\max} = \pm 0.6$ rad ($\pm 34°$)
    \item Air density: $\rho = 1.225$ kg/m$^3$
\end{itemize}

\section{Perception Layer}

\subsection{Sensor Suite}

The system relies exclusively on sensor data (never ground truth):

\paragraph{GNSS (GPS):}
Provides noisy position estimates with:
\begin{align}
\tilde{x} &= x + \mathcal{N}(0, \sigma_{\text{GNSS}}^2) \\
\tilde{y} &= y + \mathcal{N}(0, \sigma_{\text{GNSS}}^2)
\end{align}
where $\sigma_{\text{GNSS}} = 2$ m (realistic GPS accuracy).

\paragraph{IMU (Inertial Measurement Unit):}
Measures angular velocity and estimates heading:
\begin{align}
\tilde{\omega} &= \omega + \mathcal{N}(0, \sigma_{\omega}^2) \\
\tilde{\theta}_{k+1} &= \tilde{\theta}_k + \tilde{\omega} \Delta t
\end{align}

\paragraph{Wheel Encoders:}
Measures forward velocity with noise:
\begin{equation}
\tilde{v} = v + \mathcal{N}(0, \sigma_v^2)
\end{equation}

\paragraph{Camera:}
Detects lane boundaries in vehicle coordinates using computer vision. Returns left and right lane points: $\{(x_i^L, y_i^L)\}$ and $\{(x_i^R, y_i^R)\}$.

\section{Control Layer}

\subsection{Linear State-Feedback Control}

The steering controller implements linear feedback based on trajectory tracking theory.

\paragraph{Error Transformation:}
Given a goal point $(x_{la}, y_{la})$ at lookahead distance $d_{la}$, we compute errors in the robot frame:

\begin{equation}
\begin{bmatrix} e_x \\ e_y \end{bmatrix} = 
\begin{bmatrix} \cos\theta & \sin\theta \\ -\sin\theta & \cos\theta \end{bmatrix}
\begin{bmatrix} x_{la} - x \\ y_{la} - y \end{bmatrix}
\end{equation}

The heading error is:
\begin{equation}
e_\theta = \theta_{\text{goal}} - \theta
\end{equation}

\paragraph{Control Law:}
The angular velocity command is computed using linear feedback:
\begin{equation}
\omega = K_2 e_y + K_3 e_\theta
\end{equation}

where the gains are derived from pole placement:
\begin{align}
K_2 &= \frac{\omega_n^2}{|v_{\text{ref}}|} \\
K_3 &= 2 \zeta \omega_n
\end{align}

with natural frequency $\omega_n = 1.2$ rad/s and damping ratio $\zeta = 0.8$ (critically damped response).

\paragraph{Ackermann Conversion:}
Convert angular velocity to steering angle:
\begin{equation}
\delta = \arctan\left(\frac{\omega L}{v}\right)
\end{equation}

\subsection{Adaptive Lookahead}

The lookahead distance adapts based on curve state:
\begin{equation}
d_{la} = \begin{cases}
5 \text{ m} & \text{if in curve } (R < 500\text{ m}) \\
12 \text{ m} & \text{if on straight } (R \geq 500\text{ m})
\end{cases}
\end{equation}

This provides tighter control in curves and smoother tracking on straights.

\subsection{Speed Control}

\paragraph{Curve Radius Estimation:}
From detected lane points, we fit a 2nd-order polynomial:
\begin{equation}
y = ax^2 + bx + c
\end{equation}

The curvature at the midpoint is:
\begin{equation}
\kappa = \frac{|y''|}{(1 + (y')^2)^{3/2}} = \frac{|2a|}{(1 + (2ax_{\text{mid}} + b)^2)^{3/2}}
\end{equation}

The radius of curvature is:
\begin{equation}
R = \frac{1}{\kappa}
\end{equation}

\paragraph{Safe Speed Calculation:}
The maximum safe speed for a curve is derived from lateral acceleration limits:
\begin{equation}
v_{\text{safe}} = \sqrt{a_{\text{lat,max}} \cdot R} \cdot \gamma
\end{equation}

where:
\begin{itemize}
    \item $a_{\text{lat,max}} = 0.25g = 2.45$ m/s$^2$ (comfortable lateral acceleration)
    \item $\gamma = 0.90$ (10\% safety margin)
\end{itemize}

\paragraph{Curve State Machine:}
The system maintains curve state with hysteresis:
\begin{itemize}
    \item \textbf{Entry}: Immediate when $R < 500$ m
    \item \textbf{In Curve}: Track minimum safe speed: $v_{\text{min}} = \min_t v_{\text{safe}}(t)$
    \item \textbf{Exit}: After 30 consecutive frames with $R \geq 500$ m
\end{itemize}

This prevents premature speed limit release while still in the curve.

\paragraph{Brake Control:}
When overspeeding ($v > v_{\text{safe}}$), brake intensity is:
\begin{equation}
b = \text{clip}\left(0.3 + 3.0 \cdot \frac{v - v_{\text{safe}}}{v_{\text{safe}}}, \, 0.3, \, 1.0\right)
\end{equation}

providing progressive braking: 30\% for minor excess, up to 100\% for dangerous speeds.

\section{Experimental Results}

\subsection{Control Effectiveness}

The LTA system demonstrates effective lane tracking and speed management:

\paragraph{Lane Tracking Performance:}
\begin{itemize}
    \item Lateral error: $|e_y| < 0.5$ m (steady state)
    \item Heading error: $|e_\theta| < 0.1$ rad (steady state)
    \item Smooth steering with rate limiting: $|\dot{\delta}| \leq 0.05$ rad/frame
\end{itemize}

\paragraph{Speed Control Performance:}
\begin{itemize}
    \item Anticipates curves: Detects radius changes 5-12 m ahead
    \item Safe cornering: Maintains $v \leq v_{\text{safe}}$ in curves
    \item Hysteresis prevents oscillation: 10\% deadband
\end{itemize}

\subsection{Key Design Features}

\begin{enumerate}
    \item \textbf{Sensor-Only Operation}: System never uses ground truth, only noisy sensor data
    \item \textbf{Adaptive Lookahead}: 5 m in curves, 12 m on straights for optimal tracking
    \item \textbf{Minimum Speed Tracking}: Maintains lowest safe speed throughout entire curve
    \item \textbf{Exit Hysteresis}: Requires 30 consecutive straight frames before releasing speed limit
    \item \textbf{Progressive Braking}: Smooth intervention proportional to speed excess
\end{enumerate}

\section{Control Architecture Diagram}

The complete control flow:

\begin{verbatim}
┌─────────────────────────────────────────────────────────┐
│                   PERCEPTION LAYER                      │
├─────────────────────────────────────────────────────────┤
│  GNSS → (x,y)     IMU → (θ,ω)     Encoders → v         │
│  Camera → Lane Boundaries → Centerline & Curve Radius  │
└────────────────────┬────────────────────────────────────┘
                     │
                     ▼
┌─────────────────────────────────────────────────────────┐
│                    CONTROL LAYER                        │
├─────────────────────────────────────────────────────────┤
│  Curve Detection: R < 500m → in_curve = True           │
│  Adaptive Lookahead: d_la = 5m (curve) / 12m (straight)│
│                                                          │
│  STEERING CONTROL:                                      │
│    1. Find goal point at d_la on centerline            │
│    2. Compute errors: e_y, e_θ in robot frame          │
│    3. Linear feedback: ω = K2·e_y + K3·e_θ             │
│    4. Ackermann: δ = atan(ωL/v)                        │
│                                                          │
│  SPEED CONTROL:                                         │
│    1. Estimate curve radius: R = 1/κ                   │
│    2. Calculate safe speed: v_safe = √(a_lat·R)·γ     │
│    3. Track minimum in curve: v_min = min(v_safe)      │
│    4. Brake if v > v_safe: b ∝ (v-v_safe)/v_safe      │
└────────────────────┬────────────────────────────────────┘
                     │
                     ▼
┌─────────────────────────────────────────────────────────┐
│                   ACTUATION LAYER                       │
├─────────────────────────────────────────────────────────┤
│  Vehicle Dynamics:                                      │
│    ẋ = v·cos(θ)                                        │
│    ẏ = v·sin(θ)                                        │
│    θ̇ = v·tan(δ)/L                                     │
│    v̇ = (F_drive - F_drag - F_brake)/m                 │
└─────────────────────────────────────────────────────────┘
\end{verbatim}

\section{Conclusion}

The presented LTA system successfully implements autonomous lane-keeping and speed control using:
\begin{itemize}
    \item Linear state-feedback control theory from robotics fundamentals
    \item Realistic sensor modeling (GNSS, IMU, wheel encoders, camera)
    \item Physics-based vehicle dynamics (kinematics, forces, constraints)
    \item Intelligent curve detection with hysteresis
    \item Adaptive control parameters (lookahead, speed limits)
\end{itemize}

The system demonstrates effective performance in simulation, maintaining safe lane position and speeds while operating exclusively on noisy sensor data, validating the robustness of the control approach.

\section{References}

\begin{itemize}
    \item Course Slides: Linear Control III - Trajectory Tracking
    \item Course Slides: LTA Control II - Lane Tracking Applications
    \item Bicycle Model: Kinematic vehicle modeling for autonomous systems
\end{itemize}

\end{document}
